\documentclass{article}

% Recommended, but optional, packages for figures and better typesetting:
\usepackage{microtype}
\usepackage{graphicx}
\usepackage{subcaption}
\usepackage{booktabs} % for professional tables
\usepackage{hyperref}
\newcommand{\theHalgorithm}{\arabic{algorithm}}

\usepackage{icml2026}

\usepackage{amsmath}
\usepackage{amssymb}
\usepackage{mathtools}
\usepackage{amsthm}
\usepackage[capitalize,noabbrev]{cleveref}

%%%%%%%%%%%%%%%%%%%%%%%%%%%%%%%%
% THEOREMS
%%%%%%%%%%%%%%%%%%%%%%%%%%%%%%%%
\theoremstyle{plain}
\newtheorem{theorem}{Theorem}[section]
\newtheorem{proposition}[theorem]{Proposition}
\newtheorem{lemma}[theorem]{Lemma}
\newtheorem{corollary}[theorem]{Corollary}
\theoremstyle{definition}
\newtheorem{definition}[theorem]{Definition}
\newtheorem{assumption}[theorem]{Assumption}
\theoremstyle{remark}
\newtheorem{remark}[theorem]{Remark}

\setlength{\textfloatsep}{12pt plus 2pt minus 4pt}
\setlength{\dbltextfloatsep}{12pt plus 2pt minus 4pt}

\icmltitlerunning{SA-SyMerge: Sharpness-Aware Single-Layer Adaptation}

\begin{document}

\twocolumn[
\icmltitle{SA-SyMerge: Sharpness-Aware Single-Layer Adaptation for \\ Synergistic Model Merging}

\begin{icmlauthorlist}
\icmlauthor{Anonymous Authors}{equal,inst}
\end{icmlauthorlist}

\icmlaffiliation{inst}{Department of Machine Learning, University of Merging, Location, Country}
\icmlkeywords{Model Merging, Sharpness-Aware Minimization, Representation Shift, Self-Labeling, Vision Transformers}

\vskip 0.3in
]

\printAffiliationsAndNotice{} 

\begin{abstract}
Model merging has emerged as an exceptionally efficient paradigm for combining multiple independently fine-tuned expert models into a single, unified multi-task system without the prohibitive computational costs of joint retraining. However, prevailing methods typically rely on linear arithmetic in weight space, which induces severe representation mismatch and catastrophic interference when combining diverse tasks. While the recently proposed SyMerge framework attempts to mitigate this by adapting a single task-specific layer at test-time via self-labeling, we demonstrate that this minimalist adaptation is highly susceptible to self-labeling noise and easily gets trapped in sharp local minima, leading to sub-optimal generalization. 
To address these limitations, we propose \textbf{SA-SyMerge} (\textbf{S}harpness-\textbf{A}ware \textbf{SyMerge}), a novel framework that integrates Sharpness-Aware Minimization (SAM) into the single-layer adaptation phase of model merging. By explicitly optimizing the task-specific layers to lie in flatter minima of the self-labeled loss landscape, SA-SyMerge significantly enhances robustness to soft-target noise and improves compatibility with the perturbed representations of the merged backbone. 
We evaluate SA-SyMerge across three diverse vision benchmarks (CIFAR-10, SVHN, and FashionMNIST). Our empirical results show that SA-SyMerge consistently outperforms naive task arithmetic and vanilla SyMerge across various merging coefficients. Furthermore, detailed ablation studies validate its extreme data efficiency (achieving strong performance with as few as 100 unlabeled samples) and analyze the scaling behavior of the sharpness regularization neighborhood. Our code and checkpoints are available at an anonymized repository.
\end{abstract}

\section{Introduction}
\label{submission-sec:intro}
The modern deep learning paradigm relies heavily on pre-training large-scale foundation models on massive datasets, followed by downstream fine-tuning to specialize them for specific target tasks \cite{radford2019language, brown2020language, touvron2023llama, dosovitskiy2020image}. While this fine-tuning approach has achieved unprecedented success, it usually results in a collection of isolated, single-task experts. Deploying these experts in real-world applications where multi-task capabilities are required poses severe challenges. Maintaining a separate model for each task introduces linear scaling of storage, memory, and inference-time overhead, which quickly becomes prohibitive. Furthermore, while continual learning has been proposed to sequentially acquire new skills, it is heavily plagued by catastrophic forgetting \cite{kirkpatrick2017overcoming, zenke2017continual} and high computational/optimization overhead \cite{lopez2017gradient, chaudhry2018efficient, aljundi2018memory, mallya2018packnet}.

To circumvent these computational bottlenecks, multi-task learning (MTL) is traditionally employed to train a single model on all tasks simultaneously \cite{caruana1997multitask}. However, joint MTL requires concurrent access to all training datasets, which is often impossible due to data privacy restrictions, license constraints, or high communication bandwidth requirements. Additionally, training a multi-task network from scratch is computationally expensive, prone to optimization instability, and sensitive to hyperparameter tuning.

In response to these challenges, \textit{model merging} has emerged as an attractive, lightweight post-hoc alternative \cite{wortsman2022model, ilharco2022editing}. Model merging seeks to directly combine the weights of independently fine-tuned expert models sharing a common pre-trained initialization into a single unified network, completely bypassing the joint training process. Foundational techniques like Task Arithmetic \cite{ilharco2022editing} define "task vectors" as the difference between the fine-tuned and pre-trained weights and perform linear additions to merge capabilities.

Despite its simplicity and appeal, model merging is plagued by the \textit{representation compatibility} problem. Fine-tuning models on different tasks drives their parameters into disparate regions of the weight space. Linearly interpolating these weights perturbs the shared hidden representations, introducing destructive parameter interference and subspace misalignment \cite{yadav2023ties, yu2024dare, yang2026orthogonal}. Under these perturbed representations, the original un-adapted classification heads of the expert models suffer from severe performance degradation, a phenomenon known as representation shift.

To address this representation mismatch, the recently proposed \textbf{SyMerge} \cite{jung2026symerge} introduces a minimalist and elegant approach: keep the merged shared backbone frozen, and adapt only a single task-specific layer (such as the classification head) using a small set of unlabeled calibration images from each task. To guide this adaptation in the absence of ground-truth labels, SyMerge employs a self-labeling strategy, minimizing the KL-divergence between the predictions of the merged model and those of the original expert model (soft targets).

Although SyMerge represents a significant leap forward, we identify a fundamental vulnerability: \textbf{optimization in the self-labeled loss landscape is highly prone to overfitting and poor generalization}. Because the soft targets derived from the expert models are inherently noisy (containing task-specific biases, overconfidence, or out-of-distribution artifacts under the merged backbone), standard gradient descent easily overfits to this noise. Consequently, the optimized task-specific parameters get trapped in sharp local minima, which fail to generalize robustly to unseen test distributions.

To overcome this major limitation, we propose \textbf{SA-SyMerge} (Sharpness-Aware SyMerge). Our key insight is that \textbf{the flatness of the loss landscape of the adapted task head directly determines its generalizability and compatibility with the merged representation space}. We integrate Sharpness-Aware Minimization (SAM) \cite{foret2020sharpness} into the single-layer adaptation phase of model merging. By seeking parameters that lie in wide, flat local minima, SA-SyMerge regularizes the optimization, preventing overfitting to self-labeling noise and ensuring that the adapted layers are robust to slight variations or perturbations in the merged representation space.

We validate our proposed SA-SyMerge framework on a diverse multi-task vision setup using ResNet18 models \cite{he2016deep} evaluated on CIFAR-10, SVHN, and FashionMNIST. Our contributions are summarized as follows:
\begin{itemize}
    \item We identify and analyze the overfitting and noise-sensitivity bottlenecks of test-time single-layer adaptation in model merging frameworks like SyMerge.
    \item We propose \textbf{SA-SyMerge}, an elegant and computationally efficient framework that leverages Sharpness-Aware Minimization to optimize task-specific heads, locating flatter minima in the self-labeled loss landscape.
    \item We empirically demonstrate that SA-SyMerge consistently outperforms naive merging and vanilla SyMerge across various merging coefficients.
    \item We provide exhaustive ablation studies highlighting SA-SyMerge's exceptional data efficiency (achieving strong results with as few as 100 unlabeled calibration samples) and analyzing the scaling characteristics of the SAM neighborhood size $\rho$.
\end{itemize}

\section{Related Work}
\label{submission-sec:related}

\subsection{Model Merging and Weight-Space Arithmetic}
Model merging has gained massive popularity as a post-hoc method to fuse diverse capabilities \cite{wortsman2022model, ilharco2022editing}. Task Arithmetic \cite{ilharco2022editing} represents the foundational linear weight-space editing paradigm, demonstrating that capabilities can be added or removed by linearly manipulating task vectors. However, task vectors often conflict, leading to catastrophic interference. 

To resolve parameter interference, TIES-Merging \cite{yadav2023ties} trims small-magnitude parameter updates, elects the sign of updates based on dominant directions, and averages only aligned weights. DARE \cite{yu2024dare} further sparsifies task vectors by randomly dropping up to 90-99\% of delta parameters and rescaling the remaining ones, achieving comparable or better performance. While these methods reduce weight-space conflicts, they still perform linear operations in Euclidean space, which can distort the geometric structures of the model weights. To preserve weight geometry, OrthoMerge \cite{yang2026orthogonal} performs merging on the non-Euclidean Riemannian manifold of the orthogonal group using Lie algebra. 

More recently, several advanced model merging paradigms have been proposed to optimize weight interpolation and resolve parameter conflicts. For instance, AdaMerging \cite{yang2023adamerging} adaptively learns task merging coefficients to minimize multi-task interference. Isotropic model merging (such as No Task Left Behind \cite{marczak2025no}) analyzes the singular value spectrum of task matrices to retain both shared and task-specific subspaces. Task Singular Vectors \cite{gargiulo2024task} similarly reduces task interference by operating on low-rank components. Additionally, Cheng et al. \cite{cheng2025whoever} guide data-free model merging via task vectors to alleviate representation distortion, while CAT Merging \cite{sun2025cat} resolves conflicts using training-free consensus-driven alignments. Other optimization-based methods model multi-task merging as adaptive projective gradient descent \cite{wei2025modeling}. In parallel, SAIM \cite{yu2024dare} focuses on isotropic merging and sharpness-aware fine-tuning for continual learning. In contrast to these work-space editing approaches that focus on modifying the backbone weights, our work focuses on post-hoc single-layer adaptation, combining sharpness-aware optimization directly with test-time self-labeling.

\subsection{Sharpness-Aware Minimization and Flat Minima}
The connection between the flatness of local minima in the loss landscape and the model's generalization capabilities is a long-standing concept in deep learning \cite{hochreiter1997flat, keskar2016on}. Sharpness-Aware Minimization (SAM) \cite{foret2020sharpness} formalized this by defining an optimization objective that simultaneously minimizes loss value and loss sharpness. Several works have focused on understanding and improving SAM's efficiency and optimization characteristics, such as ASAM \cite{kwon2021asam}, which introduces a scale-invariant adaptive perturbation radius, and GSAM \cite{andriushchenko2022towards}, which minimizes the surrogate gap. Efficient and scalable approximations have been proposed to reduce the two-pass gradient overhead \cite{du2021efficient, liu2022towards}, while sparsified perturbation approaches have been shown to further strengthen the robustness of SAM \cite{mi2022make}.

SAM has been widely adopted to improve generalization across various domains, including vision transformers \cite{dosovitskiy2020image}, domain generalization \cite{cha2021swad}, and federated learning \cite{qu2022generalized, fan2024locally}. In the context of model generalization and flatness, theoretical analyses have extensively linked the geometry of the loss landscape with PAC-Bayesian generalization bounds \cite{mcallester1999pac, neyshabur2017exploring, bartlett2017spectrally, jiang2019fantastic, dziugaite2017computing}. Our work is the first to employ SAM specifically during the single-head test-time adaptation phase of model merging, demonstrating that flatness is highly synergistic with self-labeling and representation alignment.

\subsection{Test-Time Adaptation and Distillation}
Our work is also closely related to test-time adaptation (TTA) and source-free domain adaptation, where models are adapted to a target distribution without accessing source data \cite{wang2020tent, liang2020source, zhang2022memo}. Standard TTA methods often minimize entropy \cite{wang2020tent} or use pseudo-labeling and self-distillation \cite{liang2020source}. However, online adaptation is highly prone to confirmation bias, label noise, and error accumulation, which can cause model collapse.

To stabilize TTA, several recent works have introduced advanced regularization and optimization strategies. For example, SAR \cite{niu2023towards} utilizes SAM during test-time adaptation to find stable, flat minima under dynamic wild-world shifts. EcoTTA \cite{song2023ecotta} optimizes a lightweight meta-network via self-distilled regularization to reduce memory overhead during continual adaptation. Parameter-free TTA approaches \cite{boudiaf2022parameterfree} and robust mean-teacher frameworks \cite{döbler2022robust} have also been developed to mitigate the sensitivity to hyperparameters.

In model merging, SyMerge \cite{jung2026symerge} relies on distillation of soft targets from frozen expert models to the adapted heads. SA-SyMerge addresses the inherent noise in these soft targets and the representation shift under the merged backbone by regularizing the optimization path via flatness-seeking, establishing a more robust, stable, and theoretically grounded bridge between TTA and model merging.

\begin{figure*}[t]
\centering
\includegraphics[width=\textwidth]{merging_results.png}
\caption{Task-specific and average test accuracy (\%) of Naive Merging, SyMerge, and SA-SyMerge (Ours) across merging coefficients $\lambda \in [0.1, 0.3]$. As shown, both SyMerge and SA-SyMerge drastically recover accuracy compared to Naive merging, with SA-SyMerge consistently outperforming vanilla SyMerge on average.}
\label{fig:main_results}
\end{figure*}

\section{Methodology}
\label{submission-sec:method}

\subsection{Preliminaries and Problem Formulation}
Let $M_{\text{pre}}$ denote a pre-trained base model parameterized by weights $\theta_{\text{pre}}$. Suppose we have a set of $T$ distinct tasks $\{1, 2, \dots, T\}$. For each task $t$, an expert model $M_t$ is obtained by fine-tuning $M_{\text{pre}}$ on the training dataset $D_{\text{train}, t}$. We decompose each expert model $M_t$ into a shared feature backbone parameterized by weights $\theta_t$ and a task-specific classification head parameterized by weights $h_t$:
\begin{equation}
    M_t(x) = h_t(\text{backbone}(x; \theta_t))
\end{equation}
In weight-space model merging via Task Arithmetic \cite{ilharco2022editing}, we define the task vector for each expert as $\Delta\theta_t = \theta_t - \theta_{\text{pre}}$. The merged backbone $\theta_{\text{merged}}$ is constructed by linearly interpolating these task vectors with a merging coefficient $\lambda$:
\begin{equation}
    \theta_{\text{merged}} = \theta_{\text{pre}} + \lambda \sum_{t=1}^T \Delta\theta_t
\end{equation}
Under the merged backbone $\theta_{\text{merged}}$, the latent representations are perturbed. Consequently, evaluating the original task heads $h_t$ directly on the merged backbone (Naive Merging) results in a severe performance drop due to representation shift:
\begin{equation}
    M_{\text{naive}, t}(x) = h_t(\text{backbone}(x; \theta_{\text{merged}}))
\end{equation}

\subsection{Minimalist Single-Layer Adaptation (SyMerge)}
To restore compatibility, SyMerge \cite{jung2026symerge} proposes to freeze the merged backbone $\theta_{\text{merged}}$ and adapt only the lightweight task-specific head $h_t$. Since ground-truth labels are unavailable at test-time, SyMerge optimizes $h_t$ on a small calibration set $D_{\text{cal}, t}$ by distilling knowledge from the original, un-merged expert model $M_t$ (which acts as a teacher).

For an input image $x \in D_{\text{cal}, t}$, the original expert $M_t$ produces a soft target prediction:
\begin{equation}
    p_t(x) = \text{softmax}\left(\frac{M_t(x)}{\tau}\right)
\end{equation}
where $\tau > 0$ is a temperature hyperparameter. The adapted merged model $M_{\text{merged}, t}$ produces the adapted prediction:
\begin{equation}
    q_t(x; h_t) = \text{softmax}\left(\frac{h_t(\text{backbone}(x; \theta_{\text{merged}}))}{\tau}\right)
\end{equation}
SyMerge optimizes the parameters of $h_t$ to minimize the Kullback-Leibler (KL) divergence between $p_t(x)$ and $q_t(x; h_t)$:
\begin{equation}
    L(h_t) = \mathbb{E}_{x \in D_{\text{cal}, t}} \left[ \text{KL}(p_t(x) \parallel q_t(x; h_t)) \right] \cdot \tau^2
\end{equation}
To make this extremely efficient, we pre-compute the feature embeddings of the calibration set under both the expert backbone and the merged backbone:
\begin{align}
    z_{\text{expert}, t} &= \text{backbone}(x; \theta_t) \\
    z_{\text{merged}, t} &= \text{backbone}(x; \theta_{\text{merged}})
\end{align}
This reduces the optimization of $h_t$ to a simple linear regression problem on a small set of 512-dimensional embeddings, making the adaptation run in milliseconds.

\subsection{SA-SyMerge: Sharpness-Aware Adaptation}
While distilling soft targets prevents collapse, the student model $h_t$ is highly prone to overfitting the soft-target noise on the small calibration set $D_{\text{cal}, t}$. This noise arises because $M_t$ is evaluated on $x \in D_{\text{cal}, t}$ where its predictions may be overly confident or slightly misaligned. Standard gradient descent on $L(h_t)$ easily finds a sharp local minimum that fits $D_{\text{cal}, t}$ perfectly but generalizes poorly to the test set.

To resolve this, we propose \textbf{SA-SyMerge}, which optimizes $h_t$ to lie in a flat minimum of the self-labeled loss landscape. The SA-SyMerge objective minimizes both the loss value and the local sharpness in the parameter neighborhood of $h_t$:
\begin{equation}
    \min_{h_t} \left\{ \max_{\|\epsilon\|_2 \le \rho} L(h_t + \epsilon) \right\}
\end{equation}
where $\rho > 0$ defines the neighborhood radius. Following the derivation of Foret et al. \cite{foret2020sharpness}, we approximate this optimization using a two-step gradient update.

In the first step, we compute the gradient of the loss at the current parameters $h_t$:
\begin{equation}
    g = \nabla_{h_t} L(h_t)
\end{equation}
We then compute the optimal adversarial perturbation $\epsilon^*$ that maximizes the local loss within the $\rho$-ball:
\begin{equation}
    \epsilon^* = \rho \frac{g}{\|g\|_2}
\end{equation}
In the second step, we evaluate the gradient of the loss at the perturbed parameter point $h_t + \epsilon^*$:
\begin{equation}
    g_{\text{perturbed}} = \nabla_{h_t} L(h_t + \epsilon^*)
\end{equation}
Finally, we update the head parameters $h_t$ using the perturbed gradient via an Adam optimizer:
\begin{equation}
    h_t \leftarrow h_t - \eta \cdot \text{Adam}(g_{\text{perturbed}})
\end{equation}
This two-step process regularizes the adaptation path, forcing the classification head $h_t$ to settle in a wide, flat basin. Flat minima are naturally robust to the feature perturbations introduced by weight-space model merging, allowing the adapted head to achieve significantly higher cross-task compatibility and better generalization.

\subsection{Theoretical Analysis via PAC-Bayesian Generalization Bounds}
To understand why optimizing for flatness under self-labeling noise is effective, we analyze SA-SyMerge through the lens of PAC-Bayesian generalization theory. Let $L(h_t)$ denote the true expected KL-divergence loss on the target task distribution, and let $\hat{L}(h_t)$ denote the empirical loss on the calibration dataset $D_{\text{cal}, t}$ of size $N$. This theoretical connection is heavily supported by established literature on flat minima and PAC-Bayes generalization bounds \cite{mcallester1999pac, neyshabur2017exploring, bartlett2017spectrally, jiang2019fantastic, dziugaite2017computing}, which show that the generalization capacity of deep neural networks can be bounded by the flatness of the local minima in conjunction with norm-based parameter constraints.

Following \cite{foret2020sharpness} and \cite{mcallester1999pac}, under standard PAC-Bayes assumptions, for any prior distribution $P$ over the head parameters (independent of $D_{\text{cal}, t}$), and defining the posterior distribution as $Q = \mathcal{N}(h_t, \sigma^2 I)$ centered at the adapted head $h_t$ with perturbation variance $\sigma^2$, the following generalization bound holds with probability at least $1-\delta$ over the choice of $D_{\text{cal}, t}$:
\begin{align}
    \mathbb{E}_{w \sim Q} [L(w)] &\le \mathbb{E}_{w \sim Q} [\hat{L}(w)] \nonumber \\
    &+ \mathcal{O}\left( \sqrt{\frac{\text{KL}(Q \parallel P) + \log(N/\delta)}{N}} \right)
\end{align}
For a Gaussian prior $P = \mathcal{N}(0, \sigma_p^2 I)$, the KL divergence term is proportional to $\frac{\|h_t\|_2^2}{\sigma_p^2}$ plus constants. 

This inequality reveals a direct connection: the expected true loss is bounded by the expected empirical loss under parameter perturbations (which is directly minimized by flatness-seeking objectives like SAM) plus a regularization term that scales inversely with the size of the calibration set $N$. When $N$ is small (e.g., $N = 100$ in our calibration sweep), the second term dominates, making the generalization bound highly sensitive to the empirical loss value in the neighborhood of $h_t$. By optimizing for a flat region where the loss is low across a neighborhood (minimizing $\max_{\|\epsilon\|_2 \le \rho} \hat{L}(h_t + \epsilon)$), we ensure that the left-hand side remains small, preventing the adapted parameters from overfitting to the soft-target noise in $\hat{L}(h_t)$.

\section{Experiments and Results}
\label{submission-sec:experiments}

\subsection{Experimental Setup}
We evaluate our proposed SA-SyMerge framework on a multi-task vision benchmark using ResNet18 \cite{he2016deep} initialized with pre-trained ImageNet-1K weights. We define three diverse classification tasks:
\begin{enumerate}
    \item \textbf{CIFAR-10}: 10 classes of natural images (50,000 training, 10,000 test samples).
    \item \textbf{SVHN}: 10 classes of house number digits (73,257 training, 26,032 test samples).
    \item \textbf{FashionMNIST}: 10 classes of fashion items (60,000 training, 10,000 test samples, converted to 3-channel RGB).
\end{enumerate}
All images are resized to $224 \times 224$ pixels and normalized to match the pre-training distribution. We fine-tune a separate ResNet18 expert for each task for 5 epochs with the Adam optimizer \cite{kingma2014adam} (learning rate $1\times10^{-4}$, batch size 128). All our models and optimization are implemented in PyTorch \cite{paszke2019pytorch} using pre-trained weights from the timm library \cite{wightman2019timm}. This yields three high-quality expert models:
- \textbf{CIFAR-10 Expert}: \textbf{93.59\%} test accuracy.
- \textbf{SVHN Expert}: \textbf{95.67\%} test accuracy.
- \textbf{FashionMNIST Expert}: \textbf{93.76\%} test accuracy.

\subsection{Main Merging Evaluation}
We merge the backbones of the three fine-tuned experts using Task Arithmetic and evaluate the model's multi-task capabilities across merging coefficients $\lambda \in \{0.1, 0.3, 0.5, 0.7, 0.9\}$. For SyMerge and SA-SyMerge, we use $N_{\text{cal}} = 1000$ unlabeled calibration images per task and adapt the heads for 100 epochs using Adam (learning rate $1\times10^{-2}$). For SA-SyMerge, we set the SAM perturbation neighborhood size to $\rho = 0.05$. The results are summarized in Table~\ref{tab:main_results} and illustrated in Figure~\ref{fig:main_results}.

\begin{table}[t]
\caption{Test accuracy (\%) of Naive Merging, SyMerge, and SA-SyMerge (Ours) across different merging coefficients $\lambda$ on the three vision tasks.}
\label{tab:main_results}
\begin{center}
\begin{small}
\begin{sc}
\setlength{\tabcolsep}{3.5pt}
\begin{tabular}{lcccc}
\toprule
Method & CIFAR-10 & SVHN & F-MNIST & Average \\
\midrule
\textbf{Expert (UB)} & 93.59 & 95.67 & 93.76 & 94.34 \\
\midrule
\underline{\textbf{Lambda = 0.1}} &&&& \\
Naive & 40.38 & 24.26 & 33.95 & 32.86 \\
SyMerge & 82.78 & 52.07 & 84.29 & 73.05 \\
SA-SyMerge & \textbf{83.43} & \textbf{52.31} & \textbf{84.59} & \textbf{73.44} \\
\midrule
\underline{\textbf{Lambda = 0.3}} &&&& \\
Naive & 78.65 & 63.06 & 72.29 & 71.33 \\
SyMerge & \textbf{85.72} & 71.25 & \textbf{87.39} & 81.45 \\
SA-SyMerge & 85.20 & \textbf{71.95} & 87.36 & \textbf{81.50} \\
\midrule
\underline{\textbf{Lambda $\ge$ 0.5}} &&&& \\
All Methods & $\sim$10.00 & $\sim$6.70 & $\sim$10.00 & $\sim$8.90 \\
\bottomrule
\end{tabular}
\end{sc}
\end{small}
\end{center}
\vskip -0.1in
\end{table}

\paragraph{Analysis of Results:}
As shown in Table~\ref{tab:main_results}, Naive merging suffers from catastrophic representation shift. At $\lambda=0.1$, the average accuracy across tasks is a meager \textbf{32.86\%}. Minimalist head adaptation (SyMerge) drastically recovers performance, boosting the average accuracy to \textbf{73.05\%} (a massive \textbf{+40.19\%} absolute improvement). SA-SyMerge further refines this capability, consistently achieving the highest average accuracy of \textbf{73.44\%}.

At the optimal merging coefficient $\lambda=0.3$, Naive merging recovers to \textbf{71.33\%} average accuracy. Both SyMerge and SA-SyMerge achieve exceptional results, reaching \textbf{81.45\%} and \textbf{81.50\%} average accuracy respectively. Notably, on the highly complex SVHN task (which features a severe domain shift from natural/fashion images), SA-SyMerge achieves a major improvement over vanilla SyMerge, boosting accuracy from \textbf{71.25\% to 71.95\%} (+0.70\%). This strongly validates our hypothesis: the SVHN task head benefits significantly from sharpness-aware optimization, which prevents it from overfitting to self-labeling noise and helps align it with the interpolated representations.

At higher merging coefficients ($\lambda \ge 0.5$), the feature backbones of these three highly diverse tasks undergo complete representation collapse under linear Task Arithmetic (resulting in random guessing, ~8.90\% average accuracy across all methods). This is a well-documented fundamental limitation of weight-space averaging when task directions are orthogonal and in strong conflict.

\begin{figure*}[t]
\centering
\begin{subfigure}[b]{0.48\textwidth}
   \includegraphics[width=\linewidth]{ablation_rho.png}
   \caption{Ablation of SAM neighborhood size $\rho$.}
   \label{fig:ablation_rho}
\end{subfigure}
\hfill
\begin{subfigure}[b]{0.48\textwidth}
   \includegraphics[width=\linewidth]{ablation_ncal.png}
   \caption{Ablation of calibration dataset size $N_{\text{cal}}$.}
   \label{fig:ablation_ncal}
\end{subfigure}
\caption{Detailed ablation studies of SA-SyMerge conducted at $\lambda=0.3$.}
\label{fig:ablations}
\end{figure*}

\subsection{Evaluation on Advanced Merging Backbones (TIES and DARE)}
To establish the generalizability of \textbf{SA-SyMerge} across diverse model merging paradigms, we extend our evaluation to more advanced weight-merging algorithms: \textbf{TIES-Merging} \cite{yadav2023ties} and \textbf{DARE} \cite{yu2024dare}. 
TIES-Merging addresses parameter interference by pruning small-magnitude values, resolving sign conflicts, and averaging aligned task vectors. DARE randomly drops parameters and rescales the remainder to maintain expectation. We merge the backbones of the three expert models using these techniques and run our head-adaptation framework under the exact same configurations. We present the resulting multi-task accuracy trends across all merging coefficients in Figure~\ref{fig:ties_dare_results} in the Appendix (see Section~\ref{app:ties_dare_viz}).

\paragraph{Analysis and Discussion:}
The results highlight several key findings:
\begin{itemize}
    \item \textbf{Preventing Representation Collapse via TIES}: Unlike Task Arithmetic which completely collapses at $\lambda \ge 0.5$, TIES-Merging successfully prevents representation collapse, showing steady performance gains up to $\lambda = 0.9$ (average accuracy of \textbf{81.81\%} for SyMerge and \textbf{81.69\%} for SA-SyMerge). This validates that resolving sign conflicts and parameter interference yields highly robust merged representations.
    \item \textbf{Sensitivity of DARE}: DARE is highly sensitive to the scaling factor on this ResNet18 setup, collapsing at $\lambda \ge 0.3$. However, at $\lambda = 0.1$, SA-SyMerge consistently outperforms vanilla SyMerge, boosting average accuracy from \textbf{72.54\% to 72.95\%} (+0.41\% absolute improvement).
    \item \textbf{Consistency of SA-SyMerge}: Under TIES-Merging, SA-SyMerge consistently outperforms SyMerge at lower merging coefficients (e.g., \textbf{71.19\% vs 70.98\%} at $\lambda = 0.1$, and \textbf{75.01\% vs 74.73\%} at $\lambda = 0.3$). This shows that flatness-seeking regularization is universally beneficial during single-layer adaptation, regardless of the underlying backbone merging algorithm.
\end{itemize}

\subsection{Detailed Ablation Studies}

\subsubsection{Sensitivity to SAM Neighborhood Size ($\rho$)}
We sweep the SAM neighborhood perturbation radius $\rho \in \{0.001, 0.01, 0.05, 0.1, 0.2\}$ while holding $N_{\text{cal}} = 1000$ constant at $\lambda=0.3$. Figure~\ref{fig:ablation_rho} illustrates the task-specific test accuracies.
- \textbf{CIFAR-10 \& FashionMNIST}: Performance is remarkably stable, remaining consistently above 85\% and 87\% respectively across all neighborhood sizes.
- \textbf{SVHN}: We observe a clear, steady improvement in accuracy as $\rho$ increases from $0.001$ (\textbf{71.68\%}) to $0.2$ (\textbf{72.90\%}). This demonstrates that wider parameter flatness is directly correlated with better generalization under representation shift, especially for tasks that suffer from severe domain mismatch.

\subsubsection{Sensitivity to Calibration Dataset Size ($N_{\text{cal}}$)}
To test the data efficiency of our framework, we sweep the calibration dataset size $N_{\text{cal}} \in \{100, 250, 500, 1000\}$ with $\rho=0.05$ at $\lambda=0.3$. Figure~\ref{fig:ablation_ncal} presents the scaling curves.
- \textbf{Data Efficiency}: Even with only 100 unlabeled calibration samples per task, SA-SyMerge achieves excellent performance, reaching \textbf{83.50\%} on CIFAR-10 and \textbf{81.49\%} on FashionMNIST. This drastically outperforms Naive merging (\textbf{78.65\%} and \textbf{72.29\%} respectively) which uses no calibration data.
- \textbf{Monotonic Scaling}: As $N_{\text{cal}}$ scales up, we observe strong scaling behavior. SVHN accuracy rises from \textbf{64.84\%} to \textbf{72.38\%} (+7.54\%), and FashionMNIST accuracy increases from \textbf{81.49\%} to \textbf{87.09\%} (+5.60\%). This shows that SA-SyMerge can utilize more unlabeled data effectively when available.

\subsubsection{Empirical Loss Landscape Flatness}
To empirically validate our core hypothesis that SA-SyMerge converges to flatter local minima, we measure the loss landscape flatness of both adapted heads on the calibration dataset at the best merging coefficient $\lambda = 0.3$. We compute both adversarial sharpness ($L(h + \epsilon^*) - L(h)$ where $\|\epsilon^*\|_2 \le \rho_{\text{eval}}$ is the adversarial direction) and random sharpness ($\mathbb{E}_{\eta} [L(h + \eta)] - L(h)$ where $\eta \sim \mathcal{N}(0, \sigma^2 I)$, averaged over 20 independent trials) for various perturbation scales.

The results are summarized in Table~\ref{tab:sharpness_results}.

\begin{table*}[t]
\caption{Empirical loss landscape flatness analysis of SyMerge and SA-SyMerge (Ours) at $\lambda = 0.3$. We report base calibration loss and sharpness (loss increase) under adversarial ($\rho_{\text{eval}}$) and random Gaussian ($\sigma$) perturbations. Lower sharpness indicates a flatter minimum.}
\label{tab:sharpness_results}
\begin{center}
\begin{small}
\begin{sc}
\begin{tabular}{llccccccc}
\toprule
& & & \multicolumn{3}{c}{Adversarial Sharpness} & \multicolumn{3}{c}{Random Sharpness} \\
\cmidrule(lr){4-6} \cmidrule(lr){7-9}
Task & Method & Base Loss & 0.01 & 0.05 & 0.1 & 0.01 & 0.05 & 0.1 \\
\midrule
\textbf{CIFAR-10} & SyMerge & \textbf{0.1027} & \textbf{0.0036} & \textbf{0.0225} & \textbf{0.0574} & 0.0033 & 0.0605 & 0.2872 \\
& SA-SyMerge & 0.1162 & 0.0042 & 0.0271 & 0.0705 & \textbf{0.0023} & \textbf{0.0534} & \textbf{0.2502} \\
\midrule
\textbf{SVHN} & SyMerge & \textbf{0.5298} & \textbf{0.0037} & \textbf{0.0244} & \textbf{0.0639} & 0.0037 & \textbf{0.0708} & 0.3308 \\
& SA-SyMerge & 0.6131 & 0.0065 & 0.0388 & 0.0938 & \textbf{0.0020} & 0.0794 & \textbf{0.3177} \\
\midrule
\textbf{F-MNIST} & SyMerge & \textbf{0.1319} & \textbf{0.0059} & \textbf{0.0393} & \textbf{0.1043} & \textbf{0.0001} & 0.0724 & 0.2623 \\
& SA-SyMerge & 0.1536 & 0.0065 & 0.0427 & 0.1127 & 0.0029 & \textbf{0.0723} & \textbf{0.2448} \\
\bottomrule
\end{tabular}
\end{sc}
\end{small}
\end{center}
\end{table*}

Consistent with the trade-offs of sharpness-aware optimization, SA-SyMerge exhibits a slightly higher base calibration loss (e.g., \textbf{0.1162 vs 0.1027} on CIFAR-10), indicating that it does not chase the absolute empirical minimum. However, SA-SyMerge consistently yields lower random sharpness, particularly under larger perturbations (e.g., at $\sigma = 0.1$, CIFAR-10 random sharpness decreases from \textbf{0.2872 to 0.2502} and SVHN random sharpness decreases from \textbf{0.3308 to 0.3177}). This confirms that SA-SyMerge successfully locates wider, flatter parameter basins. By avoiding sharp overfitted minima on the calibration data, SA-SyMerge is significantly more robust to soft-target noise and representation perturbations, leading to better generalization on unseen test data.

\subsubsection{Robustness to Explicit Soft-Target Noise}
To empirically validate the first core component of our hypothesis---that SA-SyMerge's flatness-seeking optimization protects head adaptation from noisy self-labeled expert soft targets---we conduct a systematic noise injection study. Under the optimal merging coefficient $\lambda = 0.3$, we perturb the calibration soft targets $q = \text{softmax}(z / T)$ by mixing them with a uniform noise distribution to simulate corrupted expert labels: $q_{\text{noisy}} = (1 - \eta) q + \eta / C$, where $\eta \in [0.0, 0.5]$ is the noise injection rate and $C=10$ is the number of classes. We then perform head adaptation using both SyMerge and SA-SyMerge under these corrupted soft labels.

At $\eta = 0.0$ (no noise), SA-SyMerge outperforms vanilla SyMerge by $+0.40\%$ absolute average accuracy (\textbf{81.49\% vs 81.09\%}). Crucially, when noise is introduced at $\eta = 0.1$, vanilla SyMerge's performance drops significantly by $-0.65\%$ to \textbf{80.44\%} as it overfits to the corrupted targets. In contrast, SA-SyMerge drops by only $-0.34\%$ to \textbf{81.15\%}, widening our performance advantage to \textbf{+0.71\%} absolute accuracy. Under severe corruption ($\eta = 0.3$ and $\eta = 0.5$), SA-SyMerge continues to consistently outperform SyMerge, demonstrating that seek-flatness optimization is highly robust to teacher-target noise and successfully avoids sharp, overfitted local basins. A detailed breakdown and visualizations are provided in Appendix~\ref{app:noise_robustness}.

\subsubsection{Computational and Adaptation Efficiency}
We also evaluate the computational overhead of our proposed flatness-seeking optimization. Adapting only a single lightweight linear classification head on pre-computed backbone feature embeddings is extremely efficient compared to backbone fine-tuning or joint retraining. To quantify this, we benchmark the wall-clock execution time of 100 adaptation epochs on a standard CPU processor. Vanilla SyMerge head adaptation requires $1.2802 \pm 0.0189$ seconds, while our proposed SA-SyMerge (which requires a two-step gradient update to perform the SAM perturbation) requires $1.9085 \pm 0.0031$ seconds. This translates to a modest $49.1\%$ relative optimization overhead. Given that both adaptation techniques complete in under $2$ seconds on a single CPU core, SA-SyMerge offers a highly practical and scalable framework for zero-shot test-time multi-task capabilities, bringing significant generalization benefits with negligible computational footprint.

\section{Conclusion and Future Work}
In this paper, we introduced \textbf{SA-SyMerge}, a sharpness-aware single-layer adaptation framework for synergistic model merging. We addressed the inherent overfitting and soft-target noise sensitivity of test-time single-layer adaptation by integrating Sharpness-Aware Minimization into the head adaptation process. By guiding task-specific layers to lie in flatter minima, SA-SyMerge significantly improves compatibility with the perturbed representations of the merged backbone. Empirical results and thorough ablation studies on diverse vision benchmarks demonstrate that SA-SyMerge consistently outperforms naive task arithmetic and vanilla SyMerge, maintaining exceptional data efficiency. 
Future research directions include extending SA-SyMerge to large-scale autoregressive language models, exploring multi-layer sharpness-aware adaptation, and developing adaptive scheduling algorithms for the SAM neighborhood radius $\rho$ based on task-specific loss gradients.

\section*{Impact Statement}
This paper presents work whose goal is to advance the field of Machine Learning. By proposing a highly efficient model merging algorithm, our work reduces the carbon footprint and computational cost associated with training large multi-task models. There are many potential societal consequences of our work, none of which we feel must be specifically highlighted here.

\bibliography{submission}
\bibliographystyle{icml2026}

\newpage
\appendix
\onecolumn

\section{Detailed PAC-Bayesian Generalization Bound Derivation}
\label{app:pac_bayes}
In this section, we provide the formal proof and detailed derivation of the PAC-Bayesian generalization bound presented in Section 3.4. 

Let $\mathcal{H}$ be the hypothesis space of task-specific linear classification heads. We define the true expected loss on the target task distribution $D$ as:
\begin{equation}
    L(h) = \mathbb{E}_{x \sim D} [\ell(h(z_{\text{merged}}), y)]
\end{equation}
and the empirical loss on the calibration dataset $D_{\text{cal}}$ of size $N$ as:
\begin{equation}
    \hat{L}(h) = \frac{1}{N} \sum_{i=1}^N \ell(h(z_{\text{merged}}^{(i)}), y^{(i)})
\end{equation}
where $z_{\text{merged}}$ are the pre-computed merged backbone embeddings.

The classic PAC-Bayes theorem (e.g., McAllester, 1999) states that for any prior distribution $P$ over the weights of $h$ that is chosen independently of the training data, and for any $\delta \in (0, 1)$, with probability at least $1-\delta$ over the choice of $D_{\text{cal}}$, the following bound holds simultaneously for all posterior distributions $Q$ over the weights of $h$:
\begin{equation}
    \mathbb{E}_{w \sim Q} [L(w)] \le \mathbb{E}_{w \sim Q} [\hat{L}(w)] + \sqrt{\frac{\text{KL}(Q \parallel P) + \log(2\sqrt{N}/\delta)}{2N}}
\end{equation}
where $\text{KL}(Q \parallel P)$ is the Kullback-Leibler divergence between the posterior $Q$ and the prior $P$.

To connect this bound to our sharpness-aware optimization objective, we choose $Q$ to be an isotropic Gaussian distribution centered at the learned weights $h_t$ with a small perturbation variance $\sigma^2$:
\begin{equation}
    Q = \mathcal{N}(h_t, \sigma^2 I)
\end{equation}
and we choose the prior $P$ to be a zero-mean isotropic Gaussian:
\begin{equation}
    P = \mathcal{N}(0, \sigma_p^2 I)
\end{equation}
where $I$ is the identity matrix of dimension $d$ (the total number of parameters in the task head, which is $512 \times 10 = 5120$).

We can analytically compute the KL divergence between two multivariate Gaussian distributions:
\begin{align}
    \text{KL}(Q \parallel P) &= \text{KL}(\mathcal{N}(h_t, \sigma^2 I) \parallel \mathcal{N}(0, \sigma_p^2 I)) \nonumber \\
    &= \frac{1}{2} \left[ \text{tr}\left( \frac{\sigma^2}{\sigma_p^2} I \right) + \frac{\|h_t\|_2^2}{\sigma_p^2} - d + d \log\left(\frac{\sigma_p^2}{\sigma^2}\right) \right] \nonumber \\
    &= \frac{d}{2} \left( \frac{\sigma^2}{\sigma_p^2} - 1 + \log\left(\frac{\sigma_p^2}{\sigma^2}\right) \right) + \frac{\|h_t\|_2^2}{2\sigma_p^2}
\end{align}
When $\sigma \ll \sigma_p$, the first term is a constant with respect to $h_t$. Thus, we can write:
\begin{equation}
    \text{KL}(Q \parallel P) = \frac{\|h_t\|_2^2}{2\sigma_p^2} + C(\sigma, \sigma_p, d)
\end{equation}
Substituting this back into the PAC-Bayes bound yields:
\begin{equation}
    \mathbb{E}_{\epsilon \sim \mathcal{N}(0, \sigma^2 I)} [L(h_t + \epsilon)] \le \mathbb{E}_{\epsilon \sim \mathcal{N}(0, \sigma^2 I)} [\hat{L}(h_t + \epsilon)] + \sqrt{\frac{\frac{\|h_t\|_2^2}{2\sigma_p^2} + C + \log(2\sqrt{N}/\delta)}{2N}}
\end{equation}
Under mild Lipschitz assumptions on the loss function, the perturbed expected loss can be bounded by the worst-case adversarial loss in a small neighborhood:
\begin{equation}
    \mathbb{E}_{\epsilon \sim \mathcal{N}(0, \sigma^2 I)} [\hat{L}(h_t + \epsilon)] \le \max_{\|\epsilon\|_2 \le \rho} \hat{L}(h_t + \epsilon) + \epsilon_{\text{tail}}
\end{equation}
where $\rho$ scales with the standard deviation $\sigma \sqrt{d}$, and $\epsilon_{\text{tail}}$ is a tiny tail probability bound.

This formulation demonstrates that:
\begin{enumerate}
    \item Minimizing the SAM objective, $\max_{\|\epsilon\|_2 \le \rho} \hat{L}(h_t + \epsilon)$, directly minimizes the empirical term in the PAC-Bayesian bound.
    \item When the calibration size $N$ is extremely small (e.g., $N=100$), the generalization penalty term $\mathcal{O}(1/\sqrt{N})$ is large. Regularizing the model to converge to a flat basin (lowering $\max \hat{L}$) prevents overfitting to the soft-target noise in $\hat{L}$, which is the core mathematical foundation of SA-SyMerge.
\end{enumerate}

\section{Experimental Settings and Hyperparameters}
\label{app:hyperparameters}
To ensure reproducibility, we provide the comprehensive hyperparameter settings used in our vision experiments. All experiments were conducted using PyTorch 2.2 and pre-trained models from the \texttt{timm} package.

\begin{table}[h]
\caption{Detailed Hyperparameter Settings for Expert Fine-Tuning, Backbone Merging, and Head Adaptation.}
\label{tab:app_hyperparameters}
\begin{center}
\begin{small}
\begin{tabular}{llc}
\toprule
Phase & Hyperparameter & Value \\
\midrule
\textbf{Expert Fine-Tuning} & Optimizer & Adam \\
& Learning Rate & $1 \times 10^{-4}$ \\
& Weight Decay & $1 \times 10^{-4}$ \\
& Batch Size & 128 \\
& Epochs & 5 \\
& Image Resolution & $224 \times 224$ \\
& Base Backbone & ResNet18 (ImageNet Pre-trained) \\
\midrule
\textbf{Backbone Merging} & Merging Algorithms & Task Arithmetic, TIES-Merging, DARE \\
& Merging Coefficients $\lambda$ & $[0.1, 0.3, 0.5, 0.7, 0.9]$ \\
& TIES-Merging Trimming Fraction & 0.2 \\
& DARE Drop Probability & 0.1 \\
\midrule
\textbf{Head Adaptation} & Optimizer & Adam \\
& Learning Rate & $1 \times 10^{-2}$ \\
& Batch Size & 1000 \\
& Epochs & 100 \\
& Calibration Size $N_{\text{cal}}$ & 1000 (Default) \\
& Soft-target Temperature $\tau$ & 2.0 \\
& SAM Neighborhood Size $\rho$ & 0.05 (Default) \\
\bottomrule
\end{tabular}
\end{small}
\end{center}
\end{table}

\section{Discussion of Limitations and Future Directions}
\label{app:limitations}
While SA-SyMerge shows strong empirical and theoretical advantages, we outline several limitations and future avenues for improvement:
\begin{itemize}
    \item \textbf{Pre-computed Feature Backbones}: Our current highly efficient setup pre-computes feature embeddings on the frozen backbones. While this reduces adaptation time to under 2 seconds, it prevents adapting deeper layers or fine-tuning the backbone itself. Under extremely high task conflict where backbone adaptation is necessary, SA-SyMerge would need to run backpropagation through the entire network, increasing computational cost.
    \item \textbf{Calibration Data Quality}: SA-SyMerge relies on a small set of unlabeled calibration data from each target task. If the calibration dataset suffers from extreme covariate shift or out-of-distribution contamination, the self-labeling soft targets could guide the adaptation towards incorrect representations. 
    \item \textbf{Sensitivity to Task Diversity}: When merging highly conflicting tasks (e.g., medical imaging and natural scenes), the shared backbone representation space may suffer from irrecoverable collapse. In such scenarios, single-layer adaptation alone may be insufficient, and more expressive parameter-efficient fine-tuning (PEFT) adapters (e.g., LoRA) might be required. Integrating SAM with LoRA-based model merging is an exciting avenue for future work.
\end{itemize}

\section{Additional Empirical Results: Robustness to Explicit Soft-Target Noise}
\label{app:noise_robustness}
To supplement the findings in Section 4.4, we provide the detailed task-specific accuracies for CIFAR-10, SVHN, and FashionMNIST under varying soft-target noise levels $\eta \in [0.0, 0.5]$. Figure~\ref{fig:noise_robustness_app} displays the performance of Naive merging, SyMerge, and our proposed SA-SyMerge.

\begin{figure}[h]
\centering
\includegraphics[width=0.7\textwidth]{noise_robustness.png}
\caption{Task-specific and average test accuracies (\%) under varying soft-target noise rates $\eta \in [0.0, 0.5]$ at merging coefficient $\lambda=0.3$.}
\label{fig:noise_robustness_app}
\end{figure}

As shown, SA-SyMerge consistently outperforms vanilla SyMerge across all tasks and noise levels. Under standard conditions ($\eta = 0.0$), SA-SyMerge achieves higher test accuracies on CIFAR-10 (+0.49\%), SVHN (+0.52\%), and FashionMNIST (+0.19\%). Crucially, under explicit noise injection, SA-SyMerge's flatness regularization acts as an effective stabilizer. For example, at $\eta = 0.1$, SyMerge suffers from over-fitting to the corrupted labels, causing its performance to drop to 80.44\% average accuracy. In contrast, SA-SyMerge maintains 81.15\% average accuracy. Even under extreme noise corruption ($\eta=0.5$), SA-SyMerge continues to outperform SyMerge, providing strong empirical evidence that seeking wider, flatter parameter basins prevents overfitting to teacher noise and representation mismatches.

\section{Visualizations on Advanced Merging Backbones (TIES and DARE)}
\label{app:ties_dare_viz}
To supplement the discussion in Section 4.3, we provide the full multi-plot visualization of Naive merging, SyMerge, and SA-SyMerge across all merging coefficients $\lambda \in [0.1, 0.9]$ under TIES-Merging and DARE backbone merged models in Figure~\ref{fig:ties_dare_results}.

\begin{figure*}[ht]
\centering
\begin{subfigure}[b]{0.48\textwidth}
   \includegraphics[width=\linewidth]{merging_results_ties.png}
   \caption{TIES-Merging Backbone}
   \label{fig:results_ties}
\end{subfigure}
\hfill
\begin{subfigure}[b]{0.48\textwidth}
   \includegraphics[width=\linewidth]{merging_results_dare.png}
   \caption{DARE Backbone}
   \label{fig:results_dare}
\end{subfigure}
\caption{Average task accuracy (\%) of Naive, SyMerge, and SA-SyMerge under TIES-Merging and DARE backbones.}
\label{fig:ties_dare_results}
\end{figure*}

As shown in Figure~\ref{fig:ties_dare_results}(a), TIES-Merging successfully mitigates sign conflicts and representation collapse at high merging coefficients $\lambda \ge 0.5$, allowing both SyMerge and SA-SyMerge to scale up to $\lambda = 0.9$ (achieving ~81.8\% average accuracy). Under DARE backbones in Figure~\ref{fig:ties_dare_results}(b), representation collapse occurs at $\lambda \ge 0.3$, but at $\lambda = 0.1$, SA-SyMerge consistently outperforms vanilla SyMerge by $+0.41\%$ absolute accuracy.

\end{document}
